\documentclass[10pt,a4paper]{amsart}
\usepackage{amsmath,amssymb,amsthm,graphicx}
% Retained preamble for the Stacks Project chapter snapshot in this directory.
% The chapter's own preamble.tex is a copy of the shared Stacks Project preamble; the block
% below is sliced MECHANICALLY, line for line, from it. Nothing was added, removed or
% redefined; the wrapper class, the xr-hyper/hyperref wiring and the externaldocument list
% are omitted because the assembler supplies the document class.
%
% lines taken from preamble.tex: 3, 5, 9, 11, 12, 156, 157, 158, ... (51 lines total)
%
% package coverage restored from the chapter's own preamble (2026-10-13): xy, xy-2cell

\usepackage[all]{xy}
\xyoption{2cell}
\UseAllTwocells

\newenvironment{reference}{\comment}{\endcomment}
\newenvironment{slogan}{\comment}{\endcomment}
\newenvironment{history}{\comment}{\endcomment}
\theoremstyle{plain}
\newtheorem{theorem}[subsection]{Theorem}
\newtheorem{proposition}[subsection]{Proposition}
\newtheorem{lemma}[subsection]{Lemma}
\theoremstyle{definition}
\newtheorem{definition}[subsection]{Definition}
\newtheorem{example}[subsection]{Example}
\newtheorem{exercise}[subsection]{Exercise}
\newtheorem{situation}[subsection]{Situation}
\theoremstyle{remark}
\newtheorem{remark}[subsection]{Remark}
\newtheorem{remarks}[subsection]{Remarks}
\numberwithin{equation}{subsection}
\def\lim{\mathop{\mathrm{lim}}\nolimits}
\def\colim{\mathop{\mathrm{colim}}\nolimits}
\def\Spec{\mathop{\mathrm{Spec}}}
\def\Hom{\mathop{\mathrm{Hom}}\nolimits}
\def\Ext{\mathop{\mathrm{Ext}}\nolimits}
\def\SheafHom{\mathop{\mathcal{H}\!\mathit{om}}\nolimits}
\def\SheafExt{\mathop{\mathcal{E}\!\mathit{xt}}\nolimits}
\def\Sch{\mathit{Sch}}
\def\Mor{\mathop{\mathrm{Mor}}\nolimits}
\def\Ob{\mathop{\mathrm{Ob}}\nolimits}
\def\Sh{\mathop{\mathit{Sh}}\nolimits}
\def\NL{\mathop{N\!L}\nolimits}
\def\CH{\mathop{\mathrm{CH}}\nolimits}
\def\proetale{{pro\text{-}\acute{e}tale}}
\def\etale{{\acute{e}tale}}
\def\QCoh{\mathit{QCoh}}
\def\Ker{\mathop{\mathrm{Ker}}}
\def\Im{\mathop{\mathrm{Im}}}
\def\Coker{\mathop{\mathrm{Coker}}}
\def\Coim{\mathop{\mathrm{Coim}}}
\DeclareMathSymbol{\boxtimes}{\mathbin}{AMSa}{"02}
\def\QCohstack{\mathcal{QC}\!\mathit{oh}}
\def\Cohstack{\mathcal{C}\!\mathit{oh}}
\def\Spacesstack{\mathcal{S}\!\mathit{paces}}
\def\Quotfunctor{\mathrm{Quot}}
\def\Hilbfunctor{\mathrm{Hilb}}
\def\Curvesstack{\mathcal{C}\!\mathit{urves}}
\def\Polarizedstack{\mathcal{P}\!\mathit{olarized}}
\def\Complexesstack{\mathcal{C}\!\mathit{omplexes}}
\def\Pic{\mathop{\mathrm{Pic}}\nolimits}
\def\Picardstack{\mathcal{P}\!\mathit{ic}}
\def\Picardfunctor{\mathrm{Pic}}
\def\Deformationcategory{\mathcal{D}\!\mathit{ef}}
% lane-s11: discard bibliographic \begin{reference} blocks entirely (incl. any \cite inside).
\newsavebox{\limrefbox}
\ifcsname reference\endcsname\expandafter\let\csname reference\endcsname\undefined\fi\ifcsname endreference\endcsname\expandafter\let\csname endreference\endcsname\undefined\fi
\def\citation#1{\relax}
\newenvironment{reference}{\setbox\limrefbox=\vbox\bgroup\let\cite\relax}{\egroup\ignorespaces}
\usepackage{cleveref}
\crefname{originalreference}{original reference}{original references}
\Crefname{originalreference}{Original reference}{Original references}
\makeatletter
\newlabel{constructions-section-relative-glueing}{{01LG}{source}{}{sourcepointer-1}{}}
\newlabel{constructions-section-relative-glueing@cref}{{[originalreference][0][]01LG}{[1][0][]source}{}{}{}}
\makeatother
\crefname{originallemma}{original lemma}{original lemmas}
\Crefname{originallemma}{Original lemma}{Original lemmas}
\makeatletter
\newlabel{lemma-descend-finite-presentation}{{01ZM}{source}{}{sourcepointer-2}{}}
\newlabel{lemma-descend-finite-presentation@cref}{{[originallemma][0][]01ZM}{[1][0][]source}{}{}{}}
\makeatother
\crefname{originallemma}{original lemma}{original lemmas}
\Crefname{originallemma}{Original lemma}{Original lemmas}
\makeatletter
\newlabel{lemma-directed-inverse-system-has-limit}{{01YX}{source}{}{sourcepointer-3}{}}
\newlabel{lemma-directed-inverse-system-has-limit@cref}{{[originallemma][0][]01YX}{[1][0][]source}{}{}{}}
\makeatother
\title{Glueing data near a point of the limit (retrocompact variant)}
\author{The Stacks Project Authors}
\date{stacks-project:limits.tex:local-snapshot-20261012; GNU FDL 1.2 or later}
\begin{document}
\maketitle
\noindent\textit{Editorial note.} No omissions, hints or 'immediate' steps in the counted span.\par\medskip
\noindent\textit{Allowed standard prerequisites (editorial).}\begin{itemize}
\item Glueing of schemes. Schemes, glueing lemmas
\end{itemize}\medskip
\ifdefined\hypertarget\hypertarget{sourcepointer-1}{}\fi
\noindent\textit{Source reference.} Original reference 01LG: sources/stacks-chow/tags.txt, source lines 2085--2085. Cross-chapter standard prerequisite cited by the author inside the counted span; it is used as a pointer and its own proof is not part of this item.\par\medskip
\ifdefined\hypertarget\hypertarget{sourcepointer-2}{}\fi
\noindent\textit{Source reference.} Original lemma 01ZM: sources/stacks-chow/tags.txt, source lines 2599--2599. Cross-chapter standard prerequisite cited by the author inside the counted span; it is used as a pointer and its own proof is not part of this item.\par\medskip
\ifdefined\hypertarget\hypertarget{sourcepointer-3}{}\fi
\noindent\textit{Source reference.} Original lemma 01YX: sources/stacks-chow/tags.txt, source lines 2572--2572. Cross-chapter standard prerequisite cited by the author inside the counted span; it is used as a pointer and its own proof is not part of this item.\par\medskip
\begin{lemma}
\label{lemma-glueing-near-point}
Let $S$ be a scheme. Let $U \subset S$ be a retrocompact open.
Let $s \in S$ be a point in the complement of $U$. With
$V = \Spec(\mathcal{O}_{S, s}) \cap U$ there is
an equivalence of categories
$$
\colim_{s \in U' \supset U\text{ open}}
\left\{
\vcenter{
\xymatrix{
X \ar[d] \\
U'
}
}
\right\}
\longrightarrow
\left\{
\vcenter{
\xymatrix{
X' \ar[d] & Y' \ar[d] \ar[l] \ar[r] & Y \ar[d] \\
U & V \ar[l] \ar[r] & \Spec(\mathcal{O}_{S, s})
}
}
\right\}
$$
where on the left hand side the vertical arrow is of finite
presentation and on the right hand side we consider commutative diagrams
whose squares are cartesian and whose vertical arrows are
of finite presentation.
\end{lemma}

\begin{proof}
Let $W \subset S$ be an open neighbourhood of $s$. By
glueing of relative schemes, see
Constructions, Section \ref{constructions-section-relative-glueing},
the functor
$$
\left\{
\begin{matrix}
X \to U' = U \cup W \text{ of finite presentation}
\end{matrix}
\right\}
\longrightarrow
\left\{
\vcenter{
\xymatrix{
X' \ar[d] & Y' \ar[d] \ar[l] \ar[r] & Y \ar[d] \\
U & W \cap U \ar[l] \ar[r] & W
}
}
\right\}
$$
is an equivalence of categories. We have
$\mathcal{O}_{S, s} = \colim \mathcal{O}_W(W)$ where
$W$ runs over the affine open neighbourhoods of $s$.
Hence $\Spec(\mathcal{O}_{S, s}) = \lim W$ where $W$
runs over the affine open neighbourhoods of $s$.
Thus the category of schemes of finite presentation
over $\Spec(\mathcal{O}_{S, s})$ is the limit of the
category of schemes of finite presentation over
$W$ where $W$ runs over the affine open neighbourhoods
of $s$, see
Lemma \ref{lemma-descend-finite-presentation}.
For every affine open $s \in W$ we see that $U \cap W$
is quasi-compact as $U \to S$ is quasi-compact.
Hence $V = \lim W \cap U$ is a limit of
quasi-compact and quasi-separated schemes (see
Lemma \ref{lemma-directed-inverse-system-has-limit}).
Thus also the category of schemes of finite presentation
over $V$ is the limit of the
categories of schemes of finite presentation over
$W \cap U$ where $W$ runs over the affine open neighbourhoods
of $s$. The lemma follows formally from a combination
of these results.
\end{proof}

\end{document}
